\chapter{Lezione 2a -}

\section{Fisica dei Plasmi e Teoria della Dinamo: Le Collisioni Coulombiane}

Il presente documento raccoglie gli appunti relativi al Modulo 2a del corso di Fisica dei Plasmi e Teoria della Dinamo, tenuto dal Prof. Alfio Bonanno presso l'INAF – Osservatorio Astrofisico di Catania. L'argomento centrale di questa sezione è lo studio delle collisioni Coulombiane all'interno di un plasma.

\subsection{Struttura del Modulo e Bibliografia}

Il programma di questa lezione è articolato in tre punti principali:
\begin{enumerate}
\item L'analisi dello \textbf{scattering di Coulomb binario}, ovvero l'interazione a due corpi tra particelle cariche.
\item La derivazione e il significato fisico del \textbf{logaritmo di Coulomb}.
\item Lo studio delle \textbf{frequenze di collisione} e dei processi di \textbf{rilassamento} termodinamico del plasma.
\end{enumerate}

Per l'approfondimento della materia, si consiglia la consultazione dei seguenti testi: Spitzer (1962, cap. 5), Fitzpatrick (cap. 3) e Rybicki \& Lightman (cap. 5). Si farà inoltre riferimento alle note del corso (capitolo 2, paragrafi 2.1-2.3) e al Set 1 di problemi.

\section{Il Ruolo delle Collisioni in un Plasma Collettivo}

Come discusso nel Capitolo 1, al di là della lunghezza di Debye ($\lambda_D$) ogni singola carica elettrica risulta efficacemente schermata dal plasma circostante, permettendo di descrivere le interazioni a grande scala attraverso campi elettromagnetici "lisci" e continui. Tuttavia, non bisogna dimenticare che, a livello microscopico, il plasma è pur sempre costituito da un insieme discreto di cariche puntiformi. A distanze inferiori a $\lambda_D$, la natura granulare del mezzo emerge: le particelle interagiscono direttamente l'una con l'altra deflettendo le reciproche traiettorie. Questi eventi di deflessione microscopica prendono il nome di \textbf{collisioni Coulombiane}.

Le collisioni svolgono un ruolo cruciale nell'evoluzione termodinamica ed elettrodinamica del plasma. Nello specifico, l'effetto delle collisioni porta a:
\begin{itemize}
\item \textbf{Rilassare (relax)} la funzione di distribuzione delle velocità $f(\mathbf{v})$ verso uno stato di equilibrio termodinamico descritto da una distribuzione \textbf{Maxwelliana}.
\item Permettere lo \textbf{scambio di energia} termica e quantità di moto tra le diverse specie costituenti il plasma (ad esempio, tra elettroni e ioni).
\item Opporre una resistenza meccanica al flusso di corrente elettrica, originando la \textbf{resistività} macroscopica del mezzo.
\item Indurre l'emissione di radiazione da parte del plasma (come la radiazione di frenamento o bremsstrahlung, che sarà trattata in seguito).
\end{itemize}

Il messaggio concettuale fondamentale di questa lezione è che, contrariamente all'intuizione basata sulla meccanica dei gas neutri, in un plasma le dinamiche collisionali \textbf{non sono dominate dai rari urti frontali violenti (close encounters)}, bensì sono dominate dall'effetto cumulativo di un'infinità di \textbf{incontri distanti e deboli (weak, distant encounters)}.

L'intensità complessiva di queste interazioni deboli a lungo raggio è quantificata da un parametro fondamentale: il \textbf{Logaritmo di Coulomb} ($\ln \Lambda$), il cui valore numerico tipico nei plasmi si attesta tra $10$ e $40$. Un'altra caratteristica vitale che verrà dimostrata è che la frequenza (il rateo) di queste collisioni decresce drammaticamente all'aumentare della velocità relativa delle particelle, seguendo una proporzionalità del tipo $v^{-3}$.

\subsection{Implicazioni per la Teoria della Dinamo}
La comprensione delle collisioni Coulombiane non è fine a se stessa, ma costituisce il \textbf{risultato chiave} per gli sviluppi successivi del corso. Dalla teoria delle collisioni si ricaverà l'espressione per la resistività elettrica del plasma. Tale resistività, a sua volta, determina univocamente la \textbf{diffusività magnetica} $\eta$ del mezzo. Infine, la diffusività magnetica entra nella definizione del \textbf{numero di Reynolds magnetico} ($R_m = UL/\eta$), che rappresenta il parametro di controllo assoluto per l'intera Teoria della Dinamo.

\section{Il Problema a Due Corpi nello Scattering Coulombiano}

Per comprendere la dinamica delle collisioni, iniziamo analizzando l'interazione isolata tra due sole particelle cariche, definita come \textbf{problema a due corpi}. Siano $q_1$ e $q_2$ le cariche delle due particelle, con masse rispettive $m_1$ e $m_2$. Le particelle collidono con una velocità relativa iniziale $u$ e un parametro d'impatto $b$, che rappresenta la distanza minima a cui si troverebbero se non vi fosse alcuna forza deflettente tra loro.

Per semplificare l'analisi matematica di questo sistema, è opportuno effettuare un cambio di sistema di riferimento. Poiché in assenza di forze esterne il centro di massa del sistema a due particelle si muove di moto rettilineo uniforme, ci poniamo nel sistema di riferimento del centro di massa. In questo sistema, introduciamo la coordinata relativa $\mathbf{r} = \mathbf{r}_1 - \mathbf{r}_2$, che descrive la distanza vettoriale tra le due cariche.
L'equazione del moto per la coordinata relativa obbedisce alla forma:
\begin{equation}
m_r \ddot{\mathbf{r}} = \frac{q_1 q_2}{r^2} \hat{\mathbf{r}}
\end{equation}
dove $m_r$ è la \textbf{massa ridotta} del sistema, definita come:
\begin{equation}
m_r = \frac{m_1 m_2}{m_1 + m_2}
\end{equation}
Grazie a questa trasformazione, il complesso problema dell'interazione tra due corpi in movimento si riduce matematicamente al problema di \textbf{una singola particella} (avente massa $m_r$) che si muove all'interno di un \textbf{campo di potenziale Coulombiano fisso} generato nell'origine. Poiché la forza Coulombiana è centrale e conservativa, l'energia cinetica si conserva: il modulo della velocità relativa $u$ rimane invariato prima e dopo la collisione. L'unico effetto netto dell'urto è la \textbf{deflessione} della direzione del vettore velocità di un certo angolo $\theta$.

\subsection{La Distanza di Massimo Avvicinamento a Novanta Gradi ($b_{90}$)}
Per parametrizzare l'intensità della collisione, si definisce una scala spaziale naturale uguagliando l'energia potenziale elettrostatica all'energia cinetica della massa ridotta misurata a una distanza teorica. Questo bilancio energetico fornisce un \textbf{risultato chiave}, definendo il parametro d'impatto critico $b_{90}$:
\begin{equation}
b_{90} \equiv \frac{q_1 q_2}{m_r u^2}
\end{equation}
Il parametro $b_{90}$ rappresenta fisicamente la distanza di massimo avvicinamento necessaria per produrre una deflessione macroscopica e severa, convenzionalmente associata a un angolo di scattering di $90^\circ$.
In un plasma termico tipico, dove avvengono urti tra elettroni e ioni (e-i), la massa ridotta coincide praticamente con la massa dell'elettrone ($m_r \simeq m_e$). Considerando l'energia cinetica termica, questo parametro critico scala all'incirca come $b_{90} \sim Z e^2 / 3T$.

\section{Deflessione a Piccolo Angolo}

Come anticipato, in un plasma le interazioni dominanti non sono gli urti frontali (dove $b \sim b_{90}$ o minore), ma le innumerevoli deflessioni remote. Analizziamo pertanto la geometria di una collisione in cui il parametro d'impatto sia molto grande rispetto al parametro critico, ovvero $b \gg b_{90}$.

In questo regime, la deviazione subita dalla particella è talmente esigua che la traiettoria reale (\textit{path}) risulta "appena piegata" (\textit{hardly bent}). Possiamo quindi applicare l'\textbf{approssimazione impulsiva}: calcoliamo il trasferimento di quantità di moto integrando la forza Coulombiana valutata non sulla traiettoria reale (curva), ma \textbf{lungo la traiettoria rettilinea imperturbata}.
Orientando gli assi in modo che la particella incidente viaggi lungo l'asse $x$ con velocità $u$ e parametro d'impatto $b$ lungo l'asse $y$, le coordinate nel tempo sono $x = ut$ e $y = b$. La distanza dall'origine è $r = \sqrt{b^2 + u^2 t^2}$.

Scomponiamo la forza Coulombiana $F = q_1 q_2 / r^2$ nelle sue due componenti ortogonali, longitudinale ($F_{\parallel}$, parallela al moto) e trasversale ($F_{\perp}$, perpendicolare al moto):
\begin{itemize}
\item La componente \textbf{longitudinale} è $F_{\parallel} = F \frac{ut}{r} = \frac{q_1 q_2}{r^2} \frac{ut}{r}$. Questa funzione risulta essere \textbf{dispari} rispetto al tempo $t$. Fisicamente, ciò significa che la forza rallenta la particella mentre si avvicina (forza frenante) e la accelera simmetricamente mentre si allontana (forza traente). L'integrale da $-\infty$ a $+\infty$ di una funzione dispari è nullo: il trasferimento netto di quantità di moto longitudinale è zero.
\item La componente \textbf{trasversale} è $F_{\perp} = F \frac{b}{r} = \frac{q_1 q_2}{r^2} \frac{b}{r} = \frac{q_1 q_2 b}{r^3}$. Essendo una funzione pari in $t$, il suo effetto si accumula invariabilmente per tutta la durata dell'incontro. \textbf{Solo $F_{\perp}$ genera un trasferimento netto di quantità di moto}.
\end{itemize}

L'impulso trasversale netto trasferito ($\Delta p_{\perp}$) si calcola integrando la forza trasversale su tutto il tempo della collisione:
\begin{equation}
\Delta p_{\perp} = \int_{-\infty}^{\infty} F_{\perp} , dt = q_1 q_2 b \int_{-\infty}^{\infty} \frac{dt}{(b^2 + u^2 t^2)^{3/2}}
\end{equation}
.

\subsection{Soluzione dell'Integrale e Calcolo dell'Angolo di Deflessione}
Per risolvere analiticamente questo integrale, applichiamo una \textbf{sostituzione trigonometrica}. Poniamo $ut = b \tan\psi$. Calcolando il differenziale del tempo otteniamo $dt = (b/u) \sec^2\psi \, d\psi$, mentre il denominatore si trasforma in $b^2 + u^2 t^2 = b^2 \sec^2\psi$. Gli estremi di integrazione temporale da $-\infty$ a $+\infty$ si mappano negli angoli da $-\pi/2$ a $\pi/2$.

Sostituendo queste identità nell'integrale:
\begin{equation}
\int_{-\infty}^{\infty} \frac{dt}{(b^2 + u^2 t^2)^{3/2}} = \int_{-\pi/2}^{\pi/2} \frac{(b/u) \sec^2\psi}{(b^2 \sec^2\psi)^{3/2}} , d\psi = \frac{1}{ub^2} \int_{-\pi/2}^{\pi/2} \cos\psi , d\psi
\end{equation}
. L'integrale del coseno tra $-\pi/2$ e $\pi/2$ restituisce banalmente il valore $2$. Pertanto, l'impulso trasversale finale trasferito ammonta a:
\begin{equation}
\Delta p_{\perp} = \frac{2 q_1 q_2}{bu}
\end{equation}
.
Una chiave di lettura (\textit{reading}) intuitiva di questo risultato è che la particella subisce una forza media efficace d'intensità $\sim q_1 q_2 / b^2$ che agisce per una durata temporale caratteristica dell'incontro stimabile in $\sim 2b/u$.

Infine, giungiamo al \textbf{risultato chiave}: l'angolo di deflessione $\theta$. Poiché stiamo operando nel limite di piccoli angoli di scattering, l'angolo $\theta$ (espresso in radianti) si approssima al rapporto tra la variazione di quantità di moto trasversale appena calcolata ($\Delta p_{\perp}$) e la quantità di moto iniziale longitudinale ($m_r u$):
\begin{equation}
\theta \simeq \frac{\Delta p_{\perp}}{m_r u} = \frac{2 q_1 q_2}{m_r u^2 b}
\end{equation}
. Ricordando la definizione introdotta in precedenza per il parametro d'impatto critico $b_{90} = q_1 q_2 / m_r u^2$, la formula per l'angolo di deflessione si semplifica in modo elegante:
\begin{equation}
\theta \simeq \frac{2 b_{90}}{b}
\end{equation}
. Questa equazione finale evidenzia la diretta proporzionalità inversa tra l'angolo di scattering subito e il parametro d'impatto della collisione, validando a posteriori l'ipotesi di piccolo angolo per $b \gg b_{90}$.

\section{Orbita Esatta: Scattering di Rutherford}

Nelle sezioni precedenti abbiamo calcolato la deflessione con un'approssimazione valida solo per piccoli angoli. Per ricavare la traiettoria esatta senza alcuna restrizione sul parametro d'impatto, si deve risolvere l'equazione differenziale dell'orbita.
In coordinate polari $(r, \varphi)$, è prassi introdurre la variabile $s = 1/r$. Utilizzando la conservazione del momento angolare $L = m_r u b$ e l'espressione per la forza repulsiva centrale $F = q_1 q_2 s^2$, l'equazione del moto si traduce nell'\textbf{Equazione di Binet}:
\begin{equation}
\frac{d^2s}{d\varphi^2} + s = \frac{m_r F}{L^2 s^2} = \frac{q_1 q_2}{m_r u^2 b^2} = \frac{b_{90}}{b^2}
\end{equation}
Questa è un'equazione differenziale lineare del secondo ordine, non omogenea ma a coefficienti costanti.

Per risolvere il problema al contorno, assumiamo che la particella provenga dall'infinito lungo un asintoto in ingresso caratterizzato dall'angolo $\varphi = 0$. A distanza infinita ($r \to \infty$), si ha $s(0) = 0$. Inoltre, dall'approssimazione geometrica per $r \to \infty$ ($r \sin\varphi \simeq b$), si deduce la condizione sulla derivata iniziale: $s'(0) = 1/b$.
Risolvendo l'equazione con queste condizioni iniziali, si ottiene la traiettoria esatta della particella:
\begin{equation}
s(\varphi) = \frac{1}{b} \sin\varphi + \frac{b_{90}}{b^2} (\cos\varphi - 1)
\end{equation}
.

Per determinare l'angolo finale di scattering, dobbiamo calcolare l'angolo $\varphi_f$ dell'asintoto in uscita, che corrisponde nuovamente a una distanza infinita, ponendo quindi $s(\varphi_f) = 0$. Sostituendo e riordinando i termini, si ricava:
\begin{equation}
\frac{1 - \cos\varphi_f}{\sin\varphi_f} = \tan\frac{\varphi_f}{2} = \frac{b}{b_{90}}
\end{equation}
.
Dalla geometria del problema, l'angolo di deflessione effettivo $\theta$ è supplementare rispetto all'angolo asintotico, per cui $\theta = \pi - \varphi_f$. Sfruttando le identità trigonometriche, si ottiene il \textbf{risultato chiave} noto come formula per la \textbf{deflessione di Rutherford}:
\begin{equation}
\tan\frac{\theta}{2} = \frac{b_{90}}{b}
\end{equation}
.
Questa espressione esatta conferma le assunzioni precedenti: se $b = b_{90}$, si ottiene $\tan(\theta/2) = 1$, da cui $\theta = 90^\circ$; se invece $b \gg b_{90}$, la tangente approssima l'argomento e si ritrova coerentemente l'espressione per il piccolo angolo, $\theta \simeq 2b_{90}/b$.

\subsection{Visualizzazione delle Orbite Computate}
Tracciando numericamente queste orbite esatte nel piano normalizzato $(x/b_{90}, y/b_{90})$, si osserva chiaramente come la dinamica dipenda dal parametro d'impatto:
\begin{itemize}
\item Per $b = b_{90}$, la particella viene deviata esattamente di un angolo retto ($90^\circ$).
\item Per $b < b_{90}$ (urti più "centrali"), la repulsione Coulombiana vince sull'inerzia e la particella subisce un \textbf{rimbalzo all'indietro (bounces back)}.
\end{itemize}

Se confrontiamo l'angolo di deflessione esatto fornito dalla formula di Rutherford con l'approssimazione impulsiva derivata nella sezione precedente ($\theta = 2b_{90}/b$), notiamo dal grafico logaritmico che l'accordo è \textbf{eccellente per tutti i valori di $b \gtrsim 2b_{90}$}. Questa validità estesa è fondamentale perché è esattamente in questo regime di ampi parametri d'impatto che avviene la stragrande maggioranza delle collisioni all'interno di un plasma.

Nota storica: fu proprio utilizzando questa medesima formula matematica che, nel 1911, \textbf{Ernest Rutherford} spiegò i sorprendenti rimbalzi all'indietro delle particelle alfa osservati nell'esperimento di Geiger e Marsden sul foglio d'oro, giungendo così alla monumentale scoperta del nucleo atomico.

\section{La Sezione d'Urto Differenziale di Rutherford}

Per tradurre la dinamica di singola particella in un modello statistico e macroscopico, si deve calcolare la sezione d'urto differenziale ($d\sigma/d\Omega$). Consideriamo un fascio di particelle incidenti aventi parametri d'impatto compresi in un anello infinitesimo tra $b$ e $b + db$. L'area di questa corona circolare è $d\sigma = 2\pi b \, db$. A causa della diffusione di Rutherford, queste particelle verranno deflesse (scattered) in un cono solido di apertura differenziale $d\Omega = 2\pi \sin\theta \, d\theta$.

Invertendo la relazione di deflessione, $b = b_{90} \cot(\theta/2)$, calcoliamo il valore assoluto della derivata di $b$ rispetto a $\theta$:
\begin{equation}
\left| \frac{db}{d\theta} \right| = \frac{b_{90}}{2 \sin^2(\theta/2)}
\end{equation}
Sostituendo queste grandezze per trovare il rapporto $d\sigma/d\Omega$, otteniamo:
\begin{equation}
\frac{d\sigma}{d\Omega} = \frac{b}{\sin\theta} \left| \frac{db}{d\theta} \right| = \frac{b_{90} \cot(\theta/2)}{2 \sin(\theta/2) \cos(\theta/2)} \cdot \frac{b_{90}}{2 \sin^2(\theta/2)}
\end{equation}
.
Semplificando trigonometricamente i termini, perveniamo al \textbf{risultato chiave}, la celebre \textbf{sezione d'urto differenziale di Rutherford} (1911):
\begin{equation}
\frac{d\sigma}{d\Omega} = \frac{b_{90}^2}{4 \sin^4(\theta/2)}
\end{equation}
Questa espressione solleva un problema fisico e matematico enorme: la sezione d'urto \textbf{diverge con potenza $4$ ($\theta^{-4}$)} quando l'angolo di scattering tende a zero ($\theta \to 0$). L'integrazione di questa formula per ottenere la sezione d'urto totale produce pertanto un valore \textbf{infinito}.
Il significato fisico di questa divergenza asintotica è insito nella natura a \textbf{lungo raggio} della forza di Coulomb: ogni singola particella dell'universo, per quanto infinitamente distante e per quanto grande sia il suo parametro d'impatto $b$, subirà sempre e comunque una minuscola, non nulla deflessione.

\section{La Sezione d'Urto per Trasferimento di Quantità di Moto}

In ambito plasmatico non siamo primariamente interessati al calcolo degli scattering totali astratti, ma alla valutazione della quantità di moto che una particella effettivamente \textit{perde} a causa degli urti, per descrivere fenomeni macroscopici come la resistività elettrica o la forza di trascinamento (\textit{drag force}).
Durante una singola collisione con angolo di deflessione $\theta$, la componente longitudinale in avanti della velocità relativa decresce da $u$ a $u \cos\theta$. La frazione di quantità di moto relativa persa nell'urto è dunque proporzionale a $(1 - \cos\theta)$.

Definiamo la \textbf{sezione d'urto per trasferimento di quantità di moto} ($\sigma_m$) integrandola sezione d'urto differenziale pesata con questa frazione di perdita:
\begin{equation}
\sigma_m \equiv \int (1 - \cos\theta) , d\sigma = \int_0^{b_{\text{max}}} (1 - \cos\theta) 2\pi b , db
\end{equation}
(Si noti che, per anticipare il problema di divergenza, l'integrale è stato limitato a un valore massimo teorico $b_{\text{max}}$).
Per procedere col calcolo, si deve esprimere il peso $(1 - \cos\theta)$ in funzione del parametro d'impatto $b$. Usando le formule di bisezione e l'identità di Rutherford:
\begin{equation}
1 - \cos\theta = 2 \sin^2\frac{\theta}{2} = \frac{2}{1 + \cot^2(\theta/2)} = \frac{2 b_{90}^2}{b^2 + b_{90}^2}
\end{equation}
Sostituendo l'espressione pesata nell'integrale:
\begin{equation}
\sigma_m = 4\pi b_{90}^2 \int_0^{b_{\text{max}}} \frac{b , db}{b^2 + b_{90}^2} = 2\pi b_{90}^2 \ln\left( 1 + \frac{b_{\text{max}}^2}{b_{90}^2} \right)
\end{equation}
.
Assumendo che $b_{\text{max}}$ sia molto più grande del raggio di deflessione macroscopica $b_{90}$ ($b_{\text{max}} \gg b_{90}$), l'uno nell'argomento del logaritmo può essere trascurato, portando al \textbf{risultato chiave}:
\begin{equation}
\sigma_m \simeq 4\pi b_{90}^2 \ln\left(\frac{b_{\text{max}}}{b_{90}}\right)
\end{equation}
L'introduzione del fattore di peso $(1 - \cos\theta)$ ha mitigato radicalmente la grave divergenza della sezione d'urto di Rutherford ($\sim \theta^{-4}$): il trasferimento di quantità di moto rimane \textbf{divergente, ma solo logaritmicamente}. L'assegnazione di un taglio fisico finito ($b_{\text{max}}$) per sanare questo residuo di divergenza porterà alla definizione rigorosa del \textbf{Logaritmo di Coulomb}.

\section{Taglio della Divergenza (Cutting off the Divergence)}

Per calcolare esplicitamente l'integrale della sezione d'urto, è necessario risolvere il problema della divergenza asintotica introducendo dei limiti fisici di integrazione superiori e inferiori, detti \textit{cut-off}.

\subsection{Il Limite Superiore (Upper Cut-off): Lo Schermo di Debye}
Il taglio superiore si ricava analizzando la natura del campo elettrico nel plasma. Come stabilito nel Modulo 1a, a distanze superiori alla lunghezza di Debye ($\lambda_D$) il potenziale elettrostatico cessa di essere puramente Coulombiano e subisce un drastico decadimento esponenziale ($q e^{-r/\lambda_D}/r$). Oltre questa distanza critica, le forze inter-particellari si annullano e, di conseguenza, \textbf{non avviene alcuna deflessione}. Il limite superiore d'integrazione è dunque fissato a:
\begin{equation}
b_{\text{max}} = \lambda_D
\end{equation}
.

\subsection{Il Limite Inferiore (Lower Cut-off): Limite Classico vs Quantistico}
Per il limite inferiore, $b_{\text{min}}$, bisogna prestare maggiore attenzione, in quanto entrano in gioco due regimi fisici in competizione.
Il primo candidato è il limite \textbf{classico}, rappresentato dal parametro di massimo avvicinamento $b_{90}$. A distanze inferiori a $b_{90}$, gli urti comportano angoli di deflessione molto ampi (large angles), una dinamica che è già intrinsecamente contenuta e risolta nella formula esatta della sezione d'urto $\sigma_m$ derivata precedentemente, rendendo matematicamente ridondante l'integrazione a $b < b_{90}$ nell'approssimazione di piccolo angolo.

Il secondo candidato è il limite \textbf{quantistico}, imposto dal Principio di Indeterminazione di Heisenberg. Non ha senso fisico definire una traiettoria per distanze inferiori alla lunghezza d'onda di de Broglie associata alla particella. Il limite quantistico si esprime come:
\begin{equation}
b_q = \frac{\hbar}{2 m_r u}
\end{equation}
Il limite inferiore effettivo sarà quindi il maggiore tra i due valori: $b_{\text{min}} = \max(b_{90}, b_q)$.

Per determinare quale dei due effetti prevalga nell'interazione tra un elettrone e uno ione (di carica $Ze$), calcoliamo il loro rapporto:
\begin{equation}
\frac{b_{90}}{b_q} = \frac{2 Z e^2}{\hbar u} = 2 Z \frac{\alpha c}{u} \quad \text{dove} \quad \alpha = \frac{e^2}{\hbar c} \simeq \frac{1}{137}
\end{equation}
(essendo $\alpha$ la costante di struttura fine).
Da questa relazione si deduce che:
\begin{itemize}
\item Per elettroni \textbf{lenti} ($u < 2 Z \alpha c$), il rapporto è maggiore di uno, per cui domina il limite \textbf{classico} ($b_{\text{min}} = b_{90}$).
\item Per elettroni \textbf{veloci}, subentra il limite \textbf{quantistico}. Per gli elettroni termici in un plasma, l'effetto quantistico domina a temperature mediamente elevate, tipicamente per $T \gtrsim 10 Z^2 \text{ eV}$.
\end{itemize}

\section{Il Logaritmo di Coulomb Termico}

L'argomento del logaritmo individuato nell'integrale della sezione d'urto prende il nome di \textbf{Logaritmo di Coulomb} ($\ln \Lambda$). Lo deriveremo esplicitamente ponendoci nel regime del taglio \textbf{classico}.
Sostituendo l'energia cinetica della massa ridotta $m_r u^2$ con l'energia termica macroscopica $3T$, il parametro d'impatto critico diviene $b_{\text{min}} = b_{90} = Z e^2 / 3T$. Per omogeneizzare le espressioni, richiamiamo inoltre la relazione che lega la temperatura alla lunghezza di Debye elettronica: $T = 4\pi n_e e^2 \lambda_{De}^2$.

Il calcolo dell'argomento $\Lambda$ assume i seguenti passaggi formali:
\begin{equation}
\Lambda = \frac{\lambda_{De}}{b_{\text{min}}} = \frac{3 T \lambda_{De}}{Z e^2} = \frac{3 \cdot 4\pi n_e e^2 \lambda_{De}^2 \cdot \lambda_{De}}{Z e^2} = \frac{12\pi n_e \lambda_{De}^3}{Z} = \frac{9 N_D}{Z}
\end{equation}
Questa straordinaria equivalenza matematica rivela che \textbf{il parametro di plasma $N_D$ (introdotto nel Capitolo 1) riemerge naturalmente come il rapporto esatto tra il più grande e il più piccolo parametro d'impatto ammissibile}.

Nella pratica astronomica e laboratoristica, l'ente statunitense Naval Research Laboratory (NRL) fornisce delle formulazioni operative (Key result) per il logaritmo di Coulomb, valide esprimendo $n_e$ in $\text{cm}^{-3}$ e $T_e$ in $\text{eV}$:
\begin{itemize}
\item Regime classico ($T_e \lesssim 10 Z^2 \text{ eV}$): $\ln \Lambda \simeq 23 - \ln(n_e^{1/2} Z T_e^{-3/2})$.
\item Regime quantistico ($T_e \gtrsim 10 Z^2 \text{ eV}$): $\ln \Lambda \simeq 24 - \ln(n_e^{1/2} T_e^{-1})$.
\end{itemize}
Sostituendo i valori nei diversi ambienti, il logaritmo di Coulomb $\ln \Lambda$ si attesta a: $36$ nell'Intracluster medium, $24$ nel Warm ionised ISM, $18$ nella Corona Solare, $17$ in un reattore Tokamak e un anomalo $\approx 2$ nel densissimo Nucleo Solare (che rasenta la non-idealità).

\section{La Dinamica del Percorso Aleatorio (Random-Walk)}

Perché la fisica dei plasmi è dominata dagli incontri distanti e non dagli impatti violenti? La spiegazione intuitiva e matematica risiede nel quadro del \textbf{cammino aleatorio (random-walk)} nello spazio delle velocità.

Consideriamo una particella di test che attraversa un plasma a velocità $u$. Per unità di tempo, il numero di "bersagli" (particelle di campo) incrociati con parametro d'impatto compreso in un anello $db$ attorno a $b$ è pari al volume spazzato moltiplicato per la densità $n$: $2\pi b \, db \, n u$.
Ogni singola collisione deflette il vettore velocità della particella di un piccolo angolo $\theta = 2b_{90}/b$. Essendo la direzione del parametro d'impatto isotropa e casuale (random direction), la deflessione azimutale media è nulla. Tuttavia, le deviazioni quadratiche medie si \textbf{sommano in quadratura}. Il tasso di variazione temporale della varianza angolare è:
\begin{equation}
\frac{d\langle\theta^2\rangle}{dt} = n u \int_{b_{\text{min}}}^{\lambda_D} \left(\frac{2b_{90}}{b}\right)^2 2\pi b , db = 8\pi n u b_{90}^2 \int_{b_{\text{min}}}^{\lambda_D} \frac{db}{b} = 8\pi n u b_{90}^2 \ln \Lambda
\end{equation}
.
Da questo sviluppo si evincono concetti fondamentali:
\begin{itemize}
\item L'integranda per la deflessione quadratica è formata dal termine $1/b$. In scala logaritmica, questo implica che \textbf{ogni decade di parametro d'impatto contribuisce in egual misura} all'effetto complessivo (ad esempio, le collisioni tra $10$ e $100\, b_{90}$ contribuiscono tanto quanto quelle tra $100$ e $1000\, b_{90}$). È proprio l'accumulo democratico su un vastissimo range di scale spaziali che genera la forza del fattore logaritmico.
\item Definiamo la \textbf{frequenza di deflessione a $90^\circ$ accumulata} ($\nu_{90}$) come il tasso a cui la varianza angolare raggiunge l'unità ($\langle\theta^2\rangle \sim 1$). Dal calcolo appena svolto, risulta: $\nu_{90} = 8\pi n u b_{90}^2 \ln \Lambda$.
\item Se compariamo questa frequenza di rilassamento progressivo con il rateo statistico per subire un \textbf{singolo impatto catastrofico} ad ampio angolo ($n \pi b_{90}^2 u$), otteniamo un rapporto schiacciante:
\begin{equation}
\frac{\nu_{90}}{n \pi b_{90}^2 u} = 8 \ln \Lambda \approx 100 \div 300
\end{equation}
\end{itemize}
Questo dimostra matematicamente che l'effetto delle innumerevoli deviazioni infinitesimali produce una \textbf{diffusione nello spazio delle velocità} che è circa 100-300 volte più efficace dei rari urti duri. L'evoluzione cinematica di una particella in un plasma è pertanto formalizzabile con gli strumenti matematici dell'equazione di Fokker-Planck (Capitolo 7).

Questa narrazione teorica è perfettamente confermata dalle \textbf{simulazioni numeriche Monte Carlo}. Scegliendo estrazioni casuali di incontri distribuiti come $2\pi b \, db$ e imponendo micro-deflessioni $\theta = 2b_{90}/b$, i grafici statistici mostrano due andamenti limpidi:
\begin{itemize}
\item Un istogramma piatto: le varie decadi del parametro d'impatto $b/b_{90}$ cedono tutte l'esatto e identico peso (share) alla varianza totale $\langle\Delta\theta^2\rangle$, rendendo visivamente palese l'origine del $\ln \Lambda$.
\item Tracciando il "vagabondare" angolare delle singole particelle (traiettorie grigie caotiche), l'inviluppo della deflessione quadratica media (\textit{mean square deflection}, curva solida) cresce \textbf{linearmente} col tempo, secondo la legge asintotica $8\pi n u b_{90}^2 \ln \Lambda \, t$, convalidando l'interpretazione del processo come un fenomeno di pura diffusione nello spazio delle fasi.
\end{itemize}

\section{Frequenze di Collisione e Rilassamento Termodinamico}

Il calcolo della sezione d'urto per trasferimento di impulso consente di definire le scale temporali e le frequenze con cui avvengono i processi dissipativi all'interno del gas ionizzato.

\subsection{Frequenza di Collisione per un Elettrone di Test}
Consideriamo l'interazione tra una particella incidente di test (particella $a$) e un bagno di particelle di campo (particelle $b$).
Ogni singola collisione deflette la particella $a$ di un angolo $\theta$, rimuovendo una quantità di moto longitudinale pari a $m_r u (1 - \cos\theta)$. Per calcolare il \textbf{tasso totale di perdita di quantità di moto} (rate of momentum loss), integriamo questo contributo su tutte le collisioni possibili nell'unità di tempo:
\begin{equation}
\frac{dp_a}{dt} = -n_b u \int m_r u (1 - \cos\theta) , d\sigma = -m_r n_b u^2 \sigma_m \equiv -\nu_{ab} m_a u
\end{equation}
In quest'ultima uguaglianza abbiamo definito formalmente la \textbf{frequenza di collisione cinematica} $\nu_{ab}$.

Specializziamo il calcolo per il caso di un elettrone di test che si muove in un bagno di ioni di carica $Ze$. Applicando le approssimazioni appropriate: la massa ridotta è $m_r \simeq m_e$, la velocità relativa è $u = v$, il limite classico è $b_{90} = Z e^2 / m_e v^2$, e la densità di carica ionica di fondo per garantire la quasi-neutralità è $n_i Z^2 = n_e Z$.
Sostituendo questi parametri nell'espressione della frequenza definita poc'anzi, otteniamo la frequenza di collisione dipendente dalla velocità:
\begin{equation}
\nu_{ei}(v) = n_i 4\pi \left( \frac{Z e^2}{m_e v^2} \right)^2 \ln \Lambda \cdot v
\end{equation}
Raggruppando i termini si ricava il \textbf{risultato chiave}:
\begin{equation}
\nu_{ei}(v) = \frac{4\pi n_e Z e^4 \ln \Lambda}{m_e^2 v^3}
\end{equation}
.
Questa espressione analitica rivela una proprietà fisica fondamentale: la dipendenza della frequenza dalla velocità \textbf{esibisce una proporzionalità del tipo $v^{-3}$}. Questa forte dipendenza inversa emerge dal bilancio di due fattori: un fattore $v$ proporzionale al flusso di particelle in arrivo e un decadimento drastico pari a $v^{-4}$ derivante dal quadrato del raggio di deflessione $b_{90}^2$. Il significato macroscopico è evidente: \textbf{le particelle veloci collidono raramente} ("fast particles barely collide").

\subsection{Elettroni "Runaway" e Plasma Non Collisionale}
Per cogliere le implicazioni di questa dipendenza, analizziamo la distribuzione statistica delle velocità in un plasma termico. Rappresentando graficamente una distribuzione Maxwelliana $f(v)$ sovrapposta alla curva del tasso collisionale normalizzato $\nu(v)/\nu(v_{te})$, si nota che il corpo principale (bulk) della distribuzione termica giace intorno al picco di velocità probabilistica $v \approx 1.4 \, v_{te}$.

Le particelle collocate nella \textbf{coda ad alta energia} della distribuzione (the tail), per le quali ad esempio $v \gtrsim 3 \, v_{te}$, sperimentano un tasso di collisioni \textbf{da 30 a 100 volte inferiore} rispetto alle particelle del bulk. La coda veloce può considerarsi di fatto \textbf{quasi non collisionale} (nearly collisionless).
Da questa dinamica traggono origine tre conseguenze capitali per la fisica dei plasmi:
\begin{itemize}
\item L'innesco del fenomeno degli \textbf{elettroni "runaway"} (trattato nel Modulo 2b): sotto l'azione di un campo elettrico, le particelle veloci smettono di subire attrito collisionale sufficiente e accelerano indefinitamente disaccoppiandosi dal fluido.
\item Il comportamento dei \textbf{plasmi caldi} (hot plasmas), che per le alte velocità termiche in gioco risultano intrinsecamente \textbf{quasi privi di collisioni}.
\item La necessità di applicare la \textbf{correzione di Spitzer alla resistività} macroscopica del mezzo, in quanto non tutti gli elettroni trasportano corrente allo stesso modo.
\end{itemize}

\subsection{Frequenze Termiche Mediate}
Per una trattazione fluida macroscopica, occorre valutare una singola frequenza di collisione rappresentativa per l'intero ensemble termico. Ponendo approssimativamente $v \sim v_{te} = (T_e / m_e)^{1/2}$, la frequenza termica scala all'incirca come $\nu_{ei} \sim 4\pi n_e Z e^4 \ln \Lambda / (m_e^{1/2} T_e^{3/2})$.
Eseguendo l'esatta integrazione statistica su una distribuzione Maxwelliana (come discusso nel Capitolo 7), il coefficiente $4\pi$ viene rimpiazzato dal fattore $4\sqrt{2\pi}/3 \approx 3.34$. Il \textbf{risultato chiave} rigoroso per la \textbf{frequenza di collisione elettrone-ione termica} è:
\begin{equation}
\nu_{ei} = \frac{4\sqrt{2\pi}}{3} \frac{n_e Z e^4 \ln \Lambda}{m_e^{1/2} T_e^{3/2}} = 2.91 \times 10^{-6} \frac{Z n_e \ln \Lambda}{T_{e , (\text{eV})}^{3/2}} \text{ s}^{-1}
\end{equation}
.

Le proprietà di scalamento che derivano da questa formula impongono deduzioni cruciali:
\begin{itemize}
\item \textbf{Dipendenza da T:} La frequenza è proporzionale a $n T^{-3/2}$. Alzando la temperatura termodinamica del mezzo, la frequenza collisionale precipita. A $10^8 \text{ K}$, un plasma subisce un milione ($10^6$) di collisioni in meno rispetto allo stesso plasma a $10^4 \text{ K}$.
\item \textbf{Criterio di Plasma:} Valutando il rapporto adimensionale con la frequenza di plasma, $\nu_{ei} / \omega_{pe} \approx \ln \Lambda / (10 N_D)$, si nota che questo è strettamente $\ll 1$. Ciò certifica l'identità del plasma: gli elettroni compiono \textbf{moltissime oscillazioni collettive prima di incappare in una singola collisione dissipativa}.
\item \textbf{Frequenze inter-specie:} Per le particelle simili, la collisione elettrone-elettrone è dello stesso ordine di grandezza di quella elettrone-ione ($\nu_{ee} \sim \nu_{ei}$). Al contrario, le collisioni ione-ione possiedono una formula omologa con masse e temperature ioniche ($m_i, T_i$). Il rapporto scala come:
\begin{equation}
\frac{\nu_{ii}}{\nu_{ee}} \sim \left(\frac{m_e}{m_i}\right)^{1/2} \left(\frac{T_e}{T_i}\right)^{3/2} \approx \frac{1}{43}
\end{equation}
Ne consegue che gli ioni, a causa dell'inerzia maggiore, si accoppiano e rilassano molto più lentamente rispetto al gas elettronico.
\end{itemize}

\subsection{Scambio Termico e Termalizzazione (Energy Exchange)}
Mentre l'impulso può essere deviato, il \textbf{trasferimento netto di energia cinetica} (termalizzazione) per impatto tra specie di massa fortemente squilibrata è estremamente inefficiente. L'analogia cinematica è quella di palline da ping-pong (elettroni) che rimbalzano contro palle da bowling immobili (ioni).

Se una particella leggera di massa $m$ a velocità $v$ colpisce un bersaglio pesante $M \gg m$ a riposo, la velocità della particella leggera rimane virtualmente invariata in modulo (rimbalzo elastico), seppur ruotata di un angolo $\theta$. La variazione di quantità di moto al quadrato subita dal bersaglio massiccio è:
\begin{equation}
|\Delta p|^2 = |m\mathbf{v}' - m\mathbf{v}|^2 = 2m^2v^2(1 - \cos\theta)
\end{equation}
L'energia netta trasferita al bersaglio $M$ è quindi:
\begin{equation}
\Delta E = \frac{|\Delta p|^2}{2M} = \frac{m}{M} m v^2 (1 - \cos\theta) = \frac{2m}{M} E (1 - \cos\theta)
\end{equation}
.
Poiché il fattore geometrico angolare, $(1 - \cos\theta)$, è perfettamente identico a quello che pesa il calcolo per la sezione d'urto da trasferimento di impulso, si deduce per linearità il \textbf{risultato chiave} della termalizzazione: la \textbf{frequenza di scambio energetico} ($\nu_{ei}^E$) è brutalmente soppressa dal fattore di massa rispetto alla frequenza collisionale cinematica:
\begin{equation}
\nu_{ei}^E \simeq \frac{2m_e}{m_i} \nu_{ei}
\end{equation}
Questo implica fisicamente che gli elettroni collidono frequentemente per variare la propria traiettoria, ma impiegano tempi scala lunghissimi (migliaia di volte superiori) per riuscire a cedere la propria temperatura e riequilibrare termicamente gli ioni di fondo.

\section{La Gerarchia dei Tempi di Rilassamento}

Lo studio delle frequenze collisionali ci permette di comprendere la complessa termodinamica di un plasma fuori dall'equilibrio termico. Poiché lo scambio energetico dipende fortemente dal rapporto delle masse delle particelle interagenti, i diversi processi di termalizzazione e rilassamento avvengono su scale temporali nettamente distinte, instaurando una vera e propria \textbf{gerarchia del rilassamento}.

Assumendo come unità di misura temporale relativa il tempo di collisione elettrone-elettrone, la gerarchia si sviluppa nei seguenti step:
\begin{enumerate}
\item \textbf{Termalizzazione degli elettroni:} Il processo più rapido è il rilassamento del gas elettronico verso una distribuzione Maxwelliana. Il tasso caratteristico di questo processo, $\nu_{ee}$, è dell'ordine di grandezza della frequenza elettrone-ione ($\nu_{ee} \sim \nu_{ei}$). A questo evento assegniamo un \textbf{tempo relativo pari a 1}.
\item \textbf{Termalizzazione degli ioni:} Il gas ionico raggiunge a sua volta una distribuzione Maxwelliana (internalizzando una propria temperatura $T_i$) a un tasso molto inferiore, scalato per la radice del rapporto delle masse: $\nu_{ii} \sim (m_e/m_i)^{1/2} \nu_{ee}$. Per un plasma di idrogeno (dove $m_i/m_e \approx 1836$), questo richiede un \textbf{tempo relativo di $\approx 43$}.
\item \textbf{Equilibrio termico tra specie (Equalizzazione):} Il processo di gran lunga più lento è il trasferimento di energia cinetica tra elettroni e ioni al fine di pareggiare le rispettive temperature ($T_e$ e $T_i$). La frequenza di equalizzazione scala secondo il rapporto lineare delle masse: $\nu_{\text{eq}} \sim (m_e/m_i)\nu_{ei}$. Questo comporta un \textbf{tempo relativo di $\approx 1836$}.
\end{enumerate}

La dinamica dell'equalizzazione delle temperature è governata dall'equazione differenziale differenze di temperatura:
\begin{equation}
\frac{dT_e}{dt} = -\nu_{\text{eq}}(T_e - T_i)
\end{equation}
dove la frequenza di equalizzazione $\nu_{\text{eq}}$ è valutabile esplicitamente tramite la formula parametrica:
\begin{equation}
\nu_{\text{eq}} = \frac{2m_e}{m_i} \nu_{ei} = 3.2 \times 10^{-9} \frac{Z^2 n_e \ln \Lambda}{A T_{e\text{V}}^{3/2}} \text{ s}^{-1}
\end{equation}
(con $A$ indicante il numero di massa atomica dello ione).

Questa radicale lentezza nello scambio di energia inter-specie porta a una conseguenza di vitale importanza in astrofisica: la formazione di \textbf{plasmi a due temperature} (Two-temperature plasmas). Tale condizione permane in tutti gli ambienti in cui un meccanismo di riscaldamento agisce specificamente su una sola delle specie a un tasso più rapido della frequenza di riequilibrio $\nu_{\text{eq}}$. Esempi lampanti di questi regimi estremi (dove $T_i \sim 10 T_e$) si rinvengono a valle degli shock non collisionali in cui vengono riscaldati preferenzialmente gli ioni (argomento del cap. 14), nelle correnti del vento solare e nei flussi caldi di accrescimento attorno ai buchi neri.

\section{Il Cammino Libero Medio}

La frequenza di collisione ci permette di definire una fondamentale scala di lunghezza macroscopica: il \textbf{cammino libero medio} ($\lambda_{\text{mfp}}$). Esso rappresenta la distanza spaziale percorsa da una particella che viaggia alla velocità termica $v_{te}$ durante l'intervallo temporale tipico tra una collisione e l'altra ($1/\nu_{ei}$).
La derivazione si imposta combinando le grandezze note:
\begin{equation}
\lambda_{\text{mfp}} = \frac{v_{te}}{\nu_{ei}} = \left(\frac{T_e}{m_e}\right)^{1/2} \cdot \frac{3 m_e^{1/2} T_e^{3/2}}{4\sqrt{2\pi} n_e Z e^4 \ln \Lambda} = \frac{3}{4\sqrt{2\pi}} \frac{T_e^2}{n_e Z e^4 \ln \Lambda}
\end{equation}
.
Sostituendo i valori numerici delle costanti universali, perveniamo al \textbf{risultato chiave} operativo:
\begin{equation}
\lambda_{\text{mfp}} = 1.44 \times 10^{13} \frac{T_{\text{eV}}^2}{Z n_e \ln \Lambda} \text{ cm}
\end{equation}
.

L'analisi di quest'equazione rivela due peculiarità strutturali determinanti:
\begin{itemize}
\item \textbf{Cancellazione della massa:} La massa scompare totalmente dall'espressione finale. Da ciò consegue un principio di simmetria per cui gli elettroni e gli ioni, qualora si trovino alla medesima temperatura $T$, possiedono \textbf{esattamente lo stesso cammino libero medio}.
\item \textbf{Proporzionalità $\lambda_{\text{mfp}} \propto T^2/n$:} Questa dipendenza diretta dal quadrato della temperatura e inversa dalla densità implica che i plasmi astrofisici caldi e diluiti esibiscano cammini liberi medi \textbf{enormi}.
\end{itemize}

\section{Regimi del Plasma: Collisionale o Non Collisionale?}

Per stabilire il regime operativo di un sistema fisico e scegliere la teoria formale adeguata per descriverlo, è indispensabile confrontare la scala microscopica del cammino libero medio ($\lambda_{\text{mfp}}$) con la scala macroscopica, ovvero la dimensione caratteristica $L$ del sistema stesso.

Una rassegna dei principali ambienti evidenzia scenari fortemente disomogenei:
\begin{itemize}
\item \textbf{Corona Solare:} con $\lambda_{\text{mfp}} = 50 \text{ km}$ e una dimensione dei \textit{loop} coronali $L = 10^4 - 10^5 \text{ km}$, il plasma risulta pienamente \textbf{collisionale} lungo i loop magnetici.
\item \textbf{Vento Solare (a 1 AU):} con $\lambda_{\text{mfp}} \sim 0.5 \text{ AU}$ a fronte di un'estensione $L = 1 \text{ AU}$, il mezzo di propagazione si colloca in un regime \textbf{marginale}.
\item \textbf{Reattore Tokamak:} il cammino libero balza a $\lambda_{\text{mfp}} \sim 5 \text{ km}$ all'interno di una camera toroidale di raggio $L \sim 10 \text{ m}$, denotando un regime totalmente \textbf{non collisionale} lungo le linee del campo $\mathbf{B}$.
\item \textbf{Gas Intra-ammasso (Intracluster gas):} $\lambda_{\text{mfp}} \sim 10 \text{ kpc}$ e dimensione cosmologica $L \sim 1 \text{ Mpc}$, determinando un plasma che è solo \textbf{debolmente collisionale}.
\end{itemize}

L'esito di questo confronto adimensionale guida radicalmente la modellistica matematica:
\begin{itemize}
\item Quando \textbf{$\lambda_{\text{mfp}} \ll L$}: Le frequenti interazioni mantengono la distribuzione statistica $f$ costantemente Maxwelliana. Il sistema può essere legittimamente e compiutamente descritto impiegando i metodi macroscopici della teoria dei fluidi (come la Magnetoidrodinamica MHD, discussa nei cap. 9-15).
\item Quando \textbf{$\lambda_{\text{mfp}} \gtrsim L$}: L'approccio fluido crolla e si rende indispensabile l'uso formale e rigoroso della \textbf{teoria cinetica} (descritta nei cap. 5-8).
\item \textbf{Il ruolo del Campo Magnetico:} Subentra tuttavia un'avvertenza fenomenologica cruciale. Nelle dinamiche trasversali, ovvero \textit{attraverso} il campo magnetico $\mathbf{B}$, la lunghezza del passo effettivo (step length) non è più governata dal cammino libero medio $\lambda_{\text{mfp}}$, bensì dal raggio di Larmor (gyroradius). Di conseguenza, anche un plasma tecnicamente "non collisionale" può mostrare uno stringente \textbf{comportamento fluido} nella sola direzione perpendicolare al campo magnetico.
\end{itemize}

\section{Riepilogo del Modulo e Sviluppi Futuri}

Concludiamo questo secondo arco logico ricapitolando le chiavi di volta fisiche e formali dedotte nello studio delle interazioni binarie:
\begin{itemize}
\item Dalla dinamica di \textbf{Rutherford} e dalla formula esatta per la deflessione ($\tan(\theta/2) = b_{90}/b$) è derivata la sezione d'urto per il trasferimento di quantità di moto: $\sigma_m = 4\pi b_{90}^2 \ln \Lambda$.
\item Si è dimostrato matematicamente che la fisica dello scattering nel plasma è \textbf{dominata dagli incontri distanti}. Questo peso preponderante dei piccoli angoli è sintetizzato nel logaritmo di Coulomb, $\ln \Lambda = \ln(12\pi n \lambda_{De}^3)$, che assume valori tipicamente inquadrati nell'intervallo $10-40$.
\item Le dipendenze cinematiche pongono la frequenza di collisione termica proporzionale a $\nu_{ei} \propto n T^{-3/2}$ e il rateo per singola particella fortemente smorzato come $\nu(v) \propto v^{-3}$.
\item La struttura di rilassamento termodinamico inter-specie obbedisce a una severa \textbf{gerarchia temporale}, caratterizzata dal rapporto dei tempi di termalizzazione $1 : 43 : 1836$.
\end{itemize}

Nel modulo successivo (Modulo 2b), l'apparato cinetico appena ultimato verrà applicato per derivare i coefficienti di trasporto macroscopici del gas ionizzato: si andrà a quantificare formalmente la \textbf{resistività}, la \textbf{diffusività magnetica}, si esploreranno le condizioni di innesco per gli \textbf{elettroni "runaway"} e si tratterà infine l'emissione radiativa per \textbf{bremsstrahlung}.